\begin{recipe}{Kokos vis}
\kicker[Hoofdgerecht,Vis,Aziatisch,Doordeweeks]{Hoofdgerecht · vis · Aziatisch · doordeweeks}
\index[register]{Kokos vis}
\index[register]{Kokosmelk}
\index[register]{Rijst}

\meta{25 minuten \dvd Eén pan}

\lettrine{S}{nel} wokgerecht met witvis in een romige currysaus: de vis gaat
maar heel kort in de pan, de groente erna, en currypasta en kokosmelk maken
de saus af. Denk er op tijd aan: de vis moet 's ochtends al uit de vriezer.

\blockrule{Ingrediënten}
\begin{ingredients}
  \ing{200 ml}{kokosmelk}\index[register]{Kokosmelk}
  \ing{2 filets (± 300 g)}{pangasiusfilet of witvis, uit de vriezer, op tijd ontdooid}
  \ing{170 g}{langkorrelige rijst}\index[register]{Rijst}
  \ing{400 g}{Thaise, Chinese of Japanse roerbakgroente}
  \ing{70 g}{groene, gele of rode currypasta (Fairtrade Original)}
  \ing{1 el}{kokosolie, om in te bakken}
  \ing{50 g}{cashewnoten, geroosterd (evt.)}\index[register]{Cashewnoten}
  \ingb{gebakken uitjes, om te garneren}
  \ingb{sambal, voor de liefhebber}
  \ingb{komkommer, in plakjes, om erbij te serveren}
\end{ingredients}

\blockrule{Bereiding}
\tip{Bak de vis echt maar heel kort. Te lang bakken maakt de blokjes droog
en dan vallen ze uit elkaar in de pan. Let bij het uitzoeken van de
currypasta op de pittigheid: die kan sterk verschillen per merk en
soort.}
\begin{steps}
  \item Kook de rijst volgens de aanwijzingen op de verpakking, ongeveer
        10 minuten.
  \item Rooster de cashewnoten, als je ze gebruikt, zonder olie in een
        klein koekenpan op laag vuur tot ze goudbruin zijn. Zet apart.
  \item Snijd de vis in blokjes. Verhit de kokosolie goed heet in een wok of hapjespan en bak de vis
        heel kort, een paar minuten, net tot hij gaar is.
  \item Voeg de roerbakgroente toe en bak 2 tot 3 minuten mee. Voeg daarna de currypasta toe en bak een minuutje mee.
  \item Giet de kokosmelk erbij en zet het vuur laag. Laat de saus rustig
        indikken; hij mag niet te slap blijven.
  \item Serveer met de rijst, gebakken uitjes, de geroosterde cashewnoten,
        sambal voor de liefhebber en een paar plakjes komkommer erbij.
\end{steps}
\end{recipe}
